\chapter{Introduction}
\label{ch:introduction}

\section{Background and Motivation}
\label{sec:background}

Industrial control systems (ICS) operate the physical processes on which contemporary
society depends, including the electricity grid, water supply, and manufacturing
plants. The communication protocols used in these environments---IEC-104 in
substation-to-control-centre traffic, Modbus in factory floor automation,
and DNP3 in utility operations---were largely designed without security as a primary
concern, and they remain in operation today \autocite{HolikMatousekRysavyVojnar2023}.
As a result, the detection of anomalous communication is a long-running research and
engineering problem with continued operational impact.

Formal methods, and in particular automaton-based models of expected protocol
behaviour, have produced verified detectors that report whether an observed event
violates a specified property of the protocol or the network. The detector itself
returns a structured verdict that a human operator can rely on as a mathematically
grounded statement about what happened.

A separate body of work has investigated the use of large language models (LLMs) to
explain such verdicts to engineers and analysts in natural language. The motivation is
practical: operator workloads are high, the formal output is dense, and a fluent
narrative reduces the cognitive cost of acting on a verdict. The risk, however, is
that the explanation drifts from the verdict it is meant to describe. The language
model may invent claims, contradict the verdict, omit operationally critical detail,
or produce internally inconsistent text. In a safety-critical setting any of these
behaviours can lead to incorrect operational decisions.

\section{Problem Statement}
\label{sec:problem}

At the interface between a formally verified detector and an LLM-based explanation
layer there is no mechanism that prevents the explanation from making claims that the
formal verdict does not support. Existing approaches to evaluating the quality of an
explanation, such as factual-consistency scoring \autocite{Honovich2022True} and
atomic-fact decomposition \autocite{Min2023FActScore}, are valuable for offline model
selection but they remain probabilistic and applied after the fact. For safety-critical
deployments a runtime mechanism is required that uses the formal verdict as a precise
bound on what the explanation is allowed to claim, and that classifies any crossing of
that bound into a taxonomy that an operator can act on.

\section{Research Questions}
\label{sec:research-questions}

The central research question of this thesis is the following:

\begin{quote}
  Can formally verified system outputs be used as runtime bounds for detecting and
  classifying AI reasoning violations, and does such a framework produce actionable,
  operationally reliable violation detection in safety-critical AI deployments?
\end{quote}

This central question decomposes into five subordinate research questions, addressed
in subsequent chapters:

\begin{enumerate}
  \item How can the output of a formal automaton-based detection system be structured
        as a machine-readable bound against which AI reasoning can be verified at
        runtime?
  \item What taxonomy of AI reasoning violations can be formally defined with respect
        to formally verified bounds, and how do these violation types manifest in ICS
        anomaly explanation tasks?
  \item How accurately and reliably does \textsc{ReasonGuard} detect each reasoning
        violation class across different LLM architectures (local and cloud) and ICS
        event types?
  \item Under what conditions---model size, quantisation level, prompt design, event
        complexity---is AI reasoning most likely to violate formally verified bounds?
  \item What operational guidelines follow from the \textsc{ReasonGuard} evaluation for
        deploying AI explanation systems in domains where formally verified outputs are
        available?
\end{enumerate}

\section{Contributions}
\label{sec:contributions}

The thesis makes four contributions.

\paragraph{A runtime verification framework for AI reasoning.}
\textsc{ReasonGuard} is, to the best of the author's knowledge, the first runtime
verification framework that uses formally verified system outputs as bounds against
which AI explanations are checked at execution time. The architecture, the bound
schema, and the verification pipeline are presented in Chapter~\ref{ch:framework}.

\paragraph{A taxonomy of five reasoning violation classes.}
The V1--V5 taxonomy (fabricated, contradicted, over-generalised, under-specified, and
incoherent reasoning) is grounded in faithfulness research
\autocite{Maynez2020Faithfulness} and runtime verification theory
\autocite{LeuckerSchallhart2009RuntimeVerification}. Each class has a formal
definition, a detection mechanism, and an empirical incidence rate measured in this
work.

\paragraph{An empirical evaluation on ICS data.}
The framework is evaluated on the IEC-104 dataset published by the NES@FIT group at
Brno University of Technology \autocite{HolikMatousekRysavyVojnar2023}. Multiple
local and cloud language models are tested across the five canonical ICS event types,
and the violation patterns per model and per event type are reported.

\paragraph{Operational deployment guidelines.}
The empirical results are translated into deployment guidelines for organisations that
plan to use LLM-based explanation systems in domains where formally verified outputs
are available. The guidelines are mapped to the transparency and quality-management
provisions of Articles~13 and~17 of the European Union AI Act.

\section{Thesis Structure}
\label{sec:structure}

Chapter~\ref{ch:related-work} surveys the four research traditions on which this
thesis builds: formal methods and automaton-based ICS detection, LLM faithfulness and
hallucination research, runtime verification theory, and the AI transparency and
accountability requirements introduced by the EU AI Act.
Chapter~\ref{ch:framework} presents the \textsc{ReasonGuard} architecture, the bound
schema, the V1--V5 violation taxonomy, and the verification mechanism.
Chapter~\ref{ch:experimental-setup} describes the experimental setup: the datasets,
the model selection, the prompt construction, and the evaluation methodology.
Chapter~\ref{ch:results} reports the empirical results.
Chapter~\ref{ch:discussion} discusses the implications, the limitations, and the
mapping to EU AI Act compliance.
Chapter~\ref{ch:conclusion} concludes and outlines future work.
